% 中文摘要

船用螺旋桨在工作过程中易产生复杂的振动响应，其模态特性直接影响推进性能、结构强度及噪声水平。传统基于有限元方法的模态分析虽具有较高精度，但计算成本高、效率低，难以满足多参数设计空间下的快速评估需求；同时，高阶湿模态试验获取困难，现有经验公式在复杂工况下的适用性亦存在局限。针对上述问题，本文提出一种结合高精度数值建模、模态试验测试与机器学习的船用螺旋桨湿模态频率预测方法。
首先，在数值建模方面，建立了船用螺旋桨高精度有限元仿真计算方法与标准化计算规程，通过统一材料参数、边界条件、网格划分策略及湿模态求解流程，实现了干模态与湿模态频率的高一致性、高可靠性计算，为后续样本生成和频率预测奠定基础。其次，在试验测试方面，建立了高可靠性的螺旋桨模态测试方法，获得了从低阶到高阶多个模态的振型特征，并揭示了流固耦合作用下附连水效应对不同阶模态附加质量的变化规律，为湿模态仿真验证及经验公式修正提供了实验依据。在此基础上，在机器学习方面，结合参数化建模方法和自动化仿真流程，实现了多类型螺旋桨样本库的大规模扩充，并融合高精度仿真结果、试验数据及经验公式构建多源训练数据集；采用多层感知机（MLP）、XGBoost 和 LightGBM 三种机器学习算法对螺旋桨湿模态频率进行预测建模，并通过对比分析评估不同模型的预测精度与泛化能力。
研究结果表明，基于机器学习的代理模型能够有效逼近高精度有限元仿真结果，在显著降低计算时间的同时保持较高预测精度。其中，梯度提升类模型在中高阶模态预测中表现出较强稳定性，神经网络模型在低阶模态预测中具有一定优势。本文所提出的方法实现了对传统经验公式的扩充与修正，形成了兼具高精度与物理可解释性的模态频率大数据模型，为船用螺旋桨振动特性的快速预测与优化设计提供了一种高效可行的技术路径。
